% =============================================================================================
% Method section, first draft. Lines starting with "% TODO" are open decisions or checks.
% Placeholder figures: \fbox boxes; the existing schematics live in ../figures/.
% =============================================================================================

\section{Method}
\label{sec:method}

Given a source image, $K$ instance masks $\{M_k\}_{k=1}^{K}$ and one local instruction per instance,
``replace the $s_k$ with $t_k$'', multi-instance editing must satisfy three requirements at once.
\emph{Binding}: every instruction acts on its own instance and nowhere else.
\emph{Disentanglement}: within a bound region, two failures remain. In \emph{source dominance
leakage}~\cite{bindedit}, the original object survives and the edit is weak or partial; in
\emph{inter-instance leakage}~\cite{bindedit,layeredit}, adjacent edits blend into each other.
\emph{Harmonization}: the edited objects must agree with the scene and with each other in lighting,
focus, perspective and contact.

Prior work enforces disentanglement by masking the attention between image tokens~\cite{mice,breakingpoint}.
This removes leakage, but it also removes the interactions through which the model makes objects
consistent with their surroundings, and the result looks pasted in. We take the opposite approach: the
joint attention of the model is left untouched, and disentanglement is obtained by steering what each
transformer block \emph{writes} into its residual stream. Every intervention is a rotation of that write,
so it changes its direction but never adds energy to the stream (Fig.~\ref{fig:overview}).

\paragraph{Binding.}
We adopt the standard regional text binding of~\cite{mice,regionalprompting}. Each local instruction is
encoded as a separate prompt segment, and the attention mask restricts every segment and the tokens of
its instance to attend to each other; image-to-image attention is fully open. Binding is not a
contribution of this work. It is the starting point of all our experiments (``binding only'' in
Tab.~\ref{tab:ablation}), and everything below is added on top of it.

\begin{figure*}[t]
  \centering
  \fbox{\parbox[c][5cm][c]{0.95\linewidth}{\centering
    PLACEHOLDER: method overview.\\
    Left: source image, instance masks, local instructions; regional binding (prompt segment
    $\leftrightarrow$ instance tokens).
    Middle: MMDiT stack with text, latent and context token streams; at every block a hook reads the
    block's residual update and rotates it on the instance tokens.
    Right: the three references used for the rotation: the context update at the same position
    (source decoupling), the neighbouring instance (contact-local disentanglement), and the
    target/source word updates (instruction rotation).}}
  \caption{\textbf{Overview.} Instructions are bound to their instances by regional text binding; the
  joint attention is otherwise left untouched. Inside every block, the residual update of each instance
  token is rotated away from what it should not write (the source object, the neighbouring instance)
  and toward the instruction, with its length preserved.}
  \label{fig:overview}
\end{figure*}

% ---------------------------------------------------------------------------------------------
\subsection{Residual Update Rotation}
\label{sec:rotation}

\paragraph{Preliminaries.}
We build on an MMDiT editing model that processes one joint token sequence
$[\mathcal{T}; \mathcal{X}; \mathcal{C}]$: text tokens $\mathcal{T}$, latent tokens $\mathcal{X}$ of the
image being generated, and context tokens $\mathcal{C}$ that encode the source image on the same spatial
grid as $\mathcal{X}$ and are never decoded. For a latent token $t \in \mathcal{X}$ we write $\bar t \in
\mathcal{C}$ for the context token at the same position, and $\Omega_k \subset \mathcal{X}$ for the latent
tokens of instance $k$. The network consists of \emph{double-stream} blocks, in which text and image tokens
have separate weights but attend jointly, followed by \emph{single-stream} blocks, in which all tokens share
weights. Sampling runs for $N$ steps with noise levels $\sigma_0 > \dots > \sigma_{N-1}$.

\paragraph{Residual updates.}
Every block $l$ is a residual update of the hidden state of each token:
\begin{equation}
  \mathbf{h}^{l+1}_t = \mathbf{h}^{l}_t + \mathbf{D}^{l}_t ,
  \label{eq:residual}
\end{equation}
where $\mathbf{D}^{l}_t$ collects the gated attention and MLP outputs of the block. The state
$\mathbf{h}^{l}_t$ accumulates everything earlier blocks have written; the update $\mathbf{D}^{l}_t$ is what
block $l$ contributes now. We never modify the state. We replace the update of instance tokens by a steered
update $\mathbf{D}'^{\,l}_t$, so the next block receives $\mathbf{h}^{l}_t + \mathbf{D}'^{\,l}_t$ and computes
its own update from the steered stream; the effect therefore propagates through the remaining blocks and,
via the predicted velocity, to later sampling steps. Text and context tokens, and latent tokens outside all
instances, are never modified. Since a block's update contains a component shared by all tokens of the frame,
we work with the centred update
\begin{equation}
  \mathbf{d}_t = \mathbf{D}_t - \bar{\mathbf{D}},
  \qquad
  \bar{\mathbf{D}} = \frac{1}{|\mathcal{X}|}\sum_{s \in \mathcal{X}} \mathbf{D}_s ,
\end{equation}
and analogously over $\mathcal{C}$ for context tokens; the shared component $\bar{\mathbf{D}}$ is left
unchanged. We write $\mathrm{n}(\mathbf{v}) = \mathbf{v}/\lVert\mathbf{v}\rVert$ and drop the block index $l$
from now on.

\paragraph{Rotation against a reference.}
All steering terms share one operation. Given a unit \emph{reference} direction $\mathbf{r}$ and the
coefficient $p = \mathbf{d}_t^{\top}\mathbf{r}$, the update splits into a reference part
$\phi(p)\,\mathbf{r}$ and a remainder,
\begin{equation}
  \mathbf{d}_t = \phi(p)\,\mathbf{r} + \big(\mathbf{d}_t - \phi(p)\,\mathbf{r}\big),
  \label{eq:split}
\end{equation}
where $\phi(p) = p$ for a two-sided split, and $\phi(p) = \max(p,0)$ or $\min(p,0)$ when only one side of the
reference is to be acted on. A steering term proposes an edit $\Delta_t$ that removes part of the reference
component, amplifies the remainder, or adds a component along $\mathbf{r}$. The edited update is then rescaled
to its original length:
\begin{equation}
  \mathbf{d}'_t = \lVert \mathbf{d}_t \rVert \; \mathrm{n}\big(\mathbf{d}_t + \Delta_t\big) .
  \label{eq:rotation}
\end{equation}
Because $\Delta_t$ lies in $\mathrm{span}(\mathbf{d}_t, \mathbf{r})$, Eq.~\eqref{eq:rotation} is a rotation
of $\mathbf{d}_t$ within that plane (within the span of $\mathbf{d}_t$ and all references when several are
used): the block writes as much as it would have, only in a different direction.
This is what distinguishes our steering from adding a vector: its size follows the block's own activity,
it vanishes where the block writes nothing, and it injects no energy into the stream. The references are not
fixed vectors: each is read from the same forward pass, per token and per block
(Secs.~\ref{sec:source}--\ref{sec:semantic}). All strengths below are multiplied by $\sigma_i$, so the
steering fades as the sample converges.
% TODO decide whether the sigma factor is stated here or in Sec. 3.5.

\begin{figure}[t]
  \centering
  \fbox{\parbox[c][4cm][c]{0.95\linewidth}{\centering
    PLACEHOLDER: one token through the blocks.\\
    (source: ../figures/residual_token.png; panels a and b only, or redraw for one column)}}
  \caption{\textbf{Steering the residual update.} The state of a token (grey) is the sum of all block
  updates (orange). We replace the update of one block by a rotated version of it; later blocks read the
  steered state and write as usual. The state itself is never overwritten.}
  \label{fig:residual}
\end{figure}

% ---------------------------------------------------------------------------------------------
\subsection{Source Decoupling}  % TODO name
\label{sec:source}

Source dominance arises because the context tokens supply the original object at exactly the positions to
be edited. Since context tokens pass through the same blocks with the same weights and modulation as the
latent tokens, the update of the context token $\bar t$ is precisely what the block writes when it reproduces
the source at that position. We take it as the reference,
$\hat{\mathbf{s}}_t = \mathrm{n}(\mathbf{d}_{\bar t})$.
For every instance token $t \in \bigcup_k \Omega_k$, the copying part of its update is the positive component
along $\hat{\mathbf{s}}_t$, and the remainder is the edit:
\begin{align}
  \mathbf{c}_t &= \max(\mathbf{d}_t^{\top}\hat{\mathbf{s}}_t,\,0)\,\hat{\mathbf{s}}_t , \\
  \Delta^{\mathrm{src}}_t &= \beta_s\,(\mathbf{d}_t - \mathbf{c}_t) .
\end{align}
After the rotation of Eq.~\eqref{eq:rotation}, the update turns away from copying while its length is
unchanged. The copying part itself is kept, so pose, layout and lighting of the original remain available to
the edit; only the departure from the source is strengthened.

The choice of the \emph{update} rather than the context \emph{state} as reference is supported by a direct
measurement. On background tokens, which should reproduce the source, latent updates align with the context
updates far more than with the context states (mean cosine $0.37$--$0.53$ vs.\ $0.22$--$0.28$ in the
single-stream blocks; Fig.~\ref{fig:diagnostics}). Copying the source is visible in what the blocks write,
not in what has accumulated.
% TODO recompute these numbers on the final setting (they come from an earlier LoMOE run).
% TODO related: SpotEdit~\cite{spotedit} (arXiv 2512.22323, CVPR 2026) finds that hidden states of non-edited
% tokens and condition-image tokens converge to the same latent subspace during denoising (they use it to
% cache/reuse condition features). Consistent with us: if states converge, a state reference mostly carries the
% shared source content (true for every token that copies, and for every neighbour), so it cannot tell copying
% from editing; the per-block update still can. Cite here as support for update-vs-state, and in the
% residual-anchor study (state vs update references, per term).

\begin{figure}[t]
  \centering
  \fbox{\parbox[c][4cm][c]{0.95\linewidth}{\centering
    PLACEHOLDER: diagnostics from the residual stream, per step and block type.\\
    (a) cosine of latent update vs.\ context update and vs.\ context state, background and instance tokens;\\
    (b) share of instance updates along the neighbour direction near contacts, with and without steering.}}
  \caption{\textbf{Leakage is visible in the residual updates.} (a) Updates of tokens that copy the source
  align with the context \emph{updates}, much less with the context \emph{states}. (b) Near contact regions,
  a large share of an instance's update points along its neighbour.}
  \label{fig:diagnostics}
\end{figure}

% ---------------------------------------------------------------------------------------------
\subsection{Contact-Local Disentanglement}  % TODO name
\label{sec:contact}

Inter-instance leakage occurs where instances are close: tokens of one object attend strongly to an adjacent
object and adopt its appearance. For each neighbour $j$ of instance $k$, the reference is the mean centred
update of its tokens,
$\mathbf{g}_j = \mathrm{n}\big(\frac{1}{|\Omega_j|}\sum_{s\in\Omega_j}\mathbf{d}_s\big)$,
\ie what the block is currently writing into $j$. Each token $t \in \Omega_k$ is weighted toward each
neighbour by its own distance $\delta_j(t)$, in tokens, to the nearest token of $\Omega_j$:
\begin{equation}
  w_{kj}(t) = \exp\!\Big(-\frac{\max(\delta_j(t) - 1,\,0)}{\tau}\Big) .
  \label{eq:contact_weights}
\end{equation}
The part of the update that follows the neighbours, and the total contact weight of the token, are
\begin{align}
  \mathbf{q}_t &= \sum_{j \neq k} w_{kj}(t)\,(\mathbf{d}_t^{\top}\mathbf{g}_j)\,\mathbf{g}_j , \\
  \rho_t &= \min\!\Big(1,\, \sum_{j \neq k} w_{kj}(t)\Big) .
\end{align}
The edit removes the neighbour-aligned part and amplifies the instance's own remainder,
\begin{equation}
  \Delta^{\mathrm{ctc}}_t = -\alpha\,\mathbf{q}_t + \beta\,\rho_t\,(\mathbf{d}_t - \mathbf{q}_t) .
\end{equation}
The split is two-sided: writing against the neighbour entangles the two instances as much as writing toward
it. The weights are \emph{local}: they depend on each token's own distance to each neighbour, so only the parts
of an object that actually touch another one are steered, and from that neighbour only
(Fig.~\ref{fig:contact}). The interior of an object and parts far from any neighbour are left to the model,
which is where its interactions with the scene, and thus harmonization, are preserved.
% TODO schedule: the final setting applies this term at steps 0-1 only. Pending the all-step run;
% if it holds, delete this note; otherwise justify the step restriction here.

\begin{figure}[t]
  \centering
  \fbox{\parbox[c][4.5cm][c]{0.95\linewidth}{\centering
    PLACEHOLDER: contact-local steering.\\
    (sources: ../figures/ortho_decomposition.png for the split, ../figures/local_tau_steering.png for
    global vs.\ local pair weights; combine into one column-width figure)}}
  \caption{\textbf{Contact-local disentanglement.} Left: the update of a token of instance $A$ is split along
  the neighbour direction $\mathbf{g}_B$; the neighbour-aligned part is removed and the remainder amplified.
  Right: with local pair weights, each token is steered only from the neighbour it is close to.}
  \label{fig:contact}
\end{figure}

% ---------------------------------------------------------------------------------------------
\subsection{Instruction Rotation}  % TODO name
\label{sec:semantic}

The two terms above tell each instance what \emph{not} to write. The instruction itself also provides a
direction. In the single-stream blocks, text and image tokens share the weights and the residual stream, so
the update a block writes into a word token lies in the same space as the updates of image tokens. Let
$W^{s}_k$ and $W^{t}_k$ be the token positions of the source word $s_k$ and the target word $t_k$ in the local
instruction of instance $k$, and $\bar{\mathbf{D}}^{s}_k$, $\bar{\mathbf{D}}^{t}_k$ the mean updates of those
tokens. We define
\begin{equation}
  \mathbf{u}_k = \mathrm{n}\big(\bar{\mathbf{D}}^{t}_k - \bar{\mathbf{D}}^{s}_k\big),
\end{equation}
how this block treats the target word differently from the source word. Everything the two words share,
being part of the same instruction and reading the same image tokens, cancels in the difference. The edit
removes the part of an instance update that leans toward the source word and pushes toward the target word
relative to the update's own size. For $t \in \Omega_k$,
\begin{equation}
  \Delta^{\mathrm{sem}}_t = \big(\lambda\,\lVert\mathbf{d}_t\rVert
     - \gamma\,\min(\mathbf{d}_t^{\top}\mathbf{u}_k,\,0)\big)\,\mathbf{u}_k ,
\end{equation}
applied in the single-stream blocks only. As for the other terms, $\mathbf{u}_k$ is recomputed in every block
and step, and the result is rescaled by Eq.~\eqref{eq:rotation}: rather than adding a fixed guidance vector,
the update is turned toward the instruction. In contrast to classifier-free guidance, which extrapolates
between two forward passes uniformly over all tokens, the direction here is obtained within a single pass,
per block and per instance.

\begin{figure}[t]
  \centering
  \fbox{\parbox[c][4cm][c]{0.95\linewidth}{\centering
    PLACEHOLDER: instruction rotation.\\
    (source: ../figures/semantic_steering.png; left: $\mathbf{u}_k$ from the word updates, right: the
    rotation of an instance update)}}
  \caption{\textbf{Instruction rotation.} In a single-stream block, the updates written into the target and
  source words define $\mathbf{u}_k$; instance updates are rotated toward it.}
  \label{fig:semantic}
\end{figure}

% ---------------------------------------------------------------------------------------------
\subsection{Composition}
\label{sec:composition}

A token of instance $k$ may receive all three edits in the same block. Each is computed from the same update
$\mathbf{d}_t$; at sampling step $i$ they are summed and applied as a single rotation,
\begin{align}
  \Delta_t &= \sigma_i\big(\Delta^{\mathrm{src}}_t + \Delta^{\mathrm{ctc}}_t + \Delta^{\mathrm{sem}}_t\big), \\
  \mathbf{D}'_t &= \lVert \mathbf{d}_t \rVert\;\mathrm{n}\big(\mathbf{d}_t + \Delta_t\big) + \bar{\mathbf{D}} ,
\end{align}
so the result does not depend on an order of the terms and the length of the update is preserved exactly.
% TODO the code currently differs: contact and source are each rescaled to |d_t| and their increments summed,
% then the semantic term is computed from that sum and rescaled to its length. Switch the code to this joint
% form and rerun the final setting, or rewrite this paragraph to match the code.
% TODO schedule (on hold): if the all-step contact run holds, add "All terms act in every block and every
% sampling step." Otherwise add the step restriction of the contact term here.
